\section{히스테리시스와 branch 선택}\label{HYS-CH5}

히스테리시스는 같은 외부 좌표에서 둘 이상의 장기상태가 가능하고
경로가 그중 하나를 선택할 때 나타난다. 단순히 충전곡선과 방전곡선에
서로 다른 부호의 sigmoid를 쓰는 것은 상태 orientation을 바꾸므로
물리적 설명이 아니다.

\subsection{고정된 상태 orientation과 branch}

\paragraph{\physid{HYS-001} 하나의 progress 정의.}
충전과 방전 모두 동일한 균형 반응과 $\xi_j$를 사용한다.
\[
 \dot Q_{\chem,j}=a_j\dot\xi_j,\qquad
 I_j=z_jF\dot n_j
 \tag{5.1}
\]
에서 부호는 처음 선언한 forward reaction과 전류 convention으로
결정한다. branch가 바뀌어도 $z_j$, $a_j$, $\xi_j$의 의미는
뒤집지 않는다.

\paragraph{\physid{HYS-002} branch별 폐쇄.}
branch $b$의 free-energy landscape, stationary target와 mobility를
\[
 G_b(\xi,T,\bm h),\qquad
 \xi_{\tar,b}(V_{\drive},T,\bm h),\qquad
 M_b(\xi,T,\bm h)\ge0
 \tag{5.2}
\]
로 둔다. 예를 들면 gradient-flow 범주에서
\[
 \dot\xi=-M_b\frac{\partial G_b^\mathrm{eff}}{\partial\xi}
 \tag{5.3}
\]
이다. 이는 가능한 열역학적 구조이지, 특정 재료의 유일한
미시 메커니즘을 뜻하지 않는다.

\paragraph{\physid{HYS-003} metastability와 전환.}
현재 branch는 국소 안정성, nucleation condition 또는 선언된
전환면
\[
 \Phi_{b\to b'}(\bm X,V_{\drive},I)=0
 \tag{5.4}
\]
을 만날 때까지 유지될 수 있다. 전환 뒤에도 $\bm X$와 전하수지는
연속이어야 하며, 불연속 jump를 허용하면 그에 따른 잠열ㆍ방출열과
전하 변화의 적분조건을 따로 제시한다.

\subsection{상세균형의 범위}

\paragraph{\physid{KIN-050} 국소 상세균형.}
하나의 고정된 landscape와 열 reservoir 안에서 미시 가역 flux는
\[
 \ln\frac{J^+}{J^-}=\frac{\mathcal A^\chem}{RT}
 \tag{5.5}
\]
와 양립해야 한다. 평형에서 $\mathcal A^\chem=0$이고 순 flux가
사라진다. 외부 구동, 서로 다른 metastable basin 또는 소거된
내부변수가 있으면 관측된 두 branch의 유효률을 서로 나눈 값이
식 (5.5)를 만족한다고 곧바로 요구할 수 없다.

\paragraph{\physid{HYS-004} 전역 상세균형 주장 금지.}
충전용 경험률과 방전용 경험률이 다르다는 사실만으로
thermodynamic affinity를 정하거나 위반을 판정하지 않는다.
그 유효률들은 서로 다른 내부상태 또는 coarse-graining을 포함할 수
있다. 동일한 state space와 전이쌍이 입증된 뒤에만 상세균형
검증이 가능하다.

\subsection{관측 가능한 gap과 loop}

\paragraph{\physid{HYS-005} 전압 gap의 분해.}
같은 $(q,T)$에서 얻은 branch 간 gap은 개념적으로
\[
 \Delta V_\mathrm{branch}
 =\Delta V_\mathrm{meta}
 +\Delta\eta_\mathrm{ct}
 +\Delta\eta_\mathrm{transport}
 +\Delta(IR)
 +\Delta V_\mathrm{mech}
 +\Delta V_\mathrm{protocol}
 \tag{5.6}
\]
로 나눈다. 우변 항의 개별 식별에는 휴지, rate sweep,
current interruption, 온도 및 응력 정보가 필요하다.

\paragraph{\physid{HYS-006} loop area의 조건.}
\[
 W_{\mathrm{loop}}=\oint V_\app\,\dd Q
 \tag{5.7}
\]
를 소산에너지로 해석하려면 경로가 같은 온도ㆍ환경에서 시작하여
모든 관련 내부상태까지 같은 값으로 닫혀야 한다. 충전과 방전의
종단 SOC만 같고 $\bm h$ 또는 $T$가 다르면 식 (5.7)은 저장에너지
변화도 포함한다.

\paragraph{\physid{HYS-007} 휴지와 reversal.}
전류가 0이 되는 순간에도 상태는 이전 값에서 연속이며,
동역학이 허용하는 경우에만 relax한다. reversal은 외부 구동의
부호 변화이지 상태의 재초기화가 아니다. 따라서 branch 전환조건을
만나지 않은 작은 reversal에서는 return-point memory 또는
minor loop의 별도 규칙이 필요하다.

\paragraph{\physid{THM-050} 히스테리시스에서의 열.}
가역열의 부호는 branch 이름이 아니라 같은 균형 반응에 대한
$\dot Q_\rev=-IT(\partial U_\eq/\partial T)_q$로 정한다.
히스테리시스 손실은 닫힌 주기에서 상태 저장 변화가 0임을 확인한
뒤 비가역열과 비교한다.

\begin{openitem}
\physid{HYS-OPEN-01}
구동률을 0으로 보낼 때 branch gap이 남는지, 무한대기에서
nucleation이 이를 제거하는지는 observation window만으로 결정되지
않는다. dwell-time series와 독립적인 구조ㆍ응력 probe가 필요하다.
\end{openitem}

\begin{identifiability}
major loop 하나는 branch landscape, 전환장벽, mobility와 수송
과전압을 고유하게 분리하지 못한다. rate-independent memory를
주장하려면 여러 rate의 극한, minor loop, reversal와 충분한
휴지 실험을 함께 제시해야 한다.
\end{identifiability}
